\section{Limitations}
\label{sec:limitations}

Three limits bound the survey's claims. First, its unit is the benchmark, not the model: the corpus
maps how the field \emph{measures} predictive embodied intelligence and deliberately says nothing
about which world models or VLA policies are strongest. A reader looking for a model ranking will not
find one here, by design. Second, the contrast axis rests on an operational definition
(Table~\ref{tab:defs}): a benchmark counts as building a VLA-versus-world-model contrast only when
both branches are run under one protocol. A more permissive reading would move some benchmarks from
model-agnostic to partial, though the qualitative picture, that explicit contrasts are rare and
counterfactual coverage is thin, is robust to the exact threshold. Third, coverage is biased towards
recent, arXiv-indexed, English-language work reachable by web search; the funnel in
Figure~\ref{fig:prisma} records only logged stages, and the two galleries include only benchmarks
whose figures could be located and provenance-checked, so absence from a gallery is not evidence of
absence of results. Finally, the advantage-of-prediction instruments discussed in
Section~\ref{sec:gap} are assessed as candidates rather than adopted; we do not rank benchmarks by any
metric that the field has not yet validated.
